\scanpage{95}
\uline{\unsure{OK|Or}} C'est vrai.
\margin{$\unsure{*}$}

b) La droite joignant $s_0$ à $s_1$ ne contient pas d'autre sommet, i.e. les 10 \deleted{vecteurs}-sommets distincts de $s_0,s_1$ sont non prop. à $s_1=(1,0,0)$.

C'est vrai.

c) La droite joignant $s_0$ à $s_3''$ ne contient pas d'autre sommet, i.e. les 10 \unsure{aut. som.} distincts de $s_0,s_3''$ sont non prop. à
\[
s_3''=(0,1+\alpha,1)
=(0,-\alpha',1)
\]
C'est vrai.

Le deuxième pt semble plus fastidieux. Mais si on est \unsure{seulement} intéressé au cas de 4 pts \inserted{distincts} dont trois forment face, on est ramené à ceci :

Aucun des 9 autres sommets distincts de $s_0,s_1,s_2$ n'est dans le plan qu'ils engendrent, i.e. n'a de coordonnée $z$ nulle. Or c'est vrai -- et on est heureux !
% archive sheet number: 94
